\section{Conclusion}
\label{a:sec:conclusion}

The constant $A(p)$ organises the quasi-stationary mean of subcritical binary
Galton--Watson branching. It is well defined, computable by the hybrid scheme of
\cref{a:sec:practical}, and identified with the basin Koenigs value
$A(p)=2\psi_{2p}(\tfrac12)$ (\cref{a:prop:AisKoenigs}). For every $r\in(0,1)$ that
conjugacy is hypertranscendental in its own variable (\cref{a:thm:main-ht})---a theorem
with no exceptional parameters in the attracting window, since differentially algebraic
Schr\"oder functions exist only over repelling fixed points. That fact explains the
failure of every elementary closed-form search, while leaving $A(p)$ fully usable for
analysis and numerics.

What the theorem does not decide is flagged with equal care (\cref{a:sec:scope}).
Transcendence---indeed, even irrationality---of individual values $A(p_0)$, and the
elementarity or differential algebraicity of the map $p\mapsto A(p)$ in $p$, remain open:
the former because the 1993 Becker--Bergweiler values theorem is confined to equal-degree
B\"ottcher pairs and does not transfer; the latter because the question appears nowhere
in the literature. Eleven rational values nonetheless resist all inverse-symbolic
identification at over $200$-digit rigour (\cref{a:subsec:pslq}). The exceptional
elementary linearisations at $r=2$ and $r=4$ (both repelling) lie outside the branching
window; their derivations, and their place inside the classification, appear in
\cref{a:app:integrable}.

For the modelling chapters that follow, the usable conclusions are the product and series
formulae, the two-sided bounds $A>\tfrac12\Ac$, the near-critical asymptotic with
logarithmic remainder, and the practical hybrid evaluation scheme.
